% Source PDF pp. 45--50. The sentence continuing from p. 44 belongs to en-07.
The stationary value is a subgroup $C$, smooth and closed, such that $C_s=\Centr_{H_s}(T_s)$ by the density theorem (Exp. IX 4.7). It remains to show that for every point $t$ of $S$, $C_t$ is the centralizer in $H_t$ of a subtorus $T_t$ (which will imply that $C_t$ is connected). But this will follow from the following more precise lemma, applied to the family $M(\ell)_t$ of subgroups of $H_t$:
\begin{lemma}{6.8}\label{XV.6.8}
Let $G$ be a connected algebraic group defined over a field $k$, $r$ an integer $>0$, $q$ an integer prime to the characteristic\nde{The word ``residual'' has been deleted.} of $k$, $M(\ell)$ ($\ell$ ranging over the powers of $q$) a coherent family of subgroups of multiplicative type of $G$, of type $(\mathbf Z/\ell\mathbf Z)^r$ (cf. 4.6), $M$ the algebraic subgroup of $G$ generated by the $M(\ell)$ (loc. cit.), $T$ the unique maximal torus of $M$ (cf. 3.4). Then one has
\[
\Centr_G(T)=\Centr_G(M)=\Centr_G(M(\ell))\qquad\text{for $\ell$ sufficiently large.}
\]
\end{lemma}
The last equality is quite clear. To prove the first, introduce the center $Z$ of $G$, $G'=G/Z$, $M'$ (resp. $T'$) the image of $M$ (resp. $T$) in $G'$, $K$ the inverse image of $T'$ in $G$ (i.e. the algebraic subgroup of $G$ generated by $T$ and $Z$). It plainly suffices to prove that $K\supset M$, hence that $T'=M'$. Now $M$ is smooth and connected (4.6) and $G'$ is affine (Exp. XII 6.1), so $M'$ is the direct product of its maximal torus $T'$ (Exp. XII 6.6 d)) and a unipotent group (BIBLE 4 Th. 4) (one may suppose $k$ algebraically closed). The image of $M(\ell)$ in $M'$ is therefore necessarily contained in $T'$. Thus the inverse image of $T'$ in $M$ contains $M(\ell)$ for every $\ell$, hence is equal to $M$. Consequently $M'=T'$. This proves 6.8 and therefore 6.7.

This being so, let us prove 6.6. We have reduced to the case where $S$ is the spectrum of a discrete valuation ring $A$, which one may moreover suppose complete with algebraically closed residue field. Replacing $A$ by its normalization in a finite extension of its field of fractions, one may suppose that $(G_t)_{\mathrm{red}}$ is smooth\nde{Details or references to be given here\ldots}. It is clear that, to prove 6.6, one may replace $G$ by the identity component of the schematic closure in $G$ of $(G_t)^0_{\mathrm{red}}$, hence suppose that $G$ is flat over $S$, with connected fibers, and that $G_t$ is smooth.

\noindent a) Suppose that $H_s$ is the centralizer in $(G_s)_{\mathrm{red}}$ of a torus $T_s$ and let us show that $H_t$ is then the centralizer in $G_t$ of a subtorus of $G_t$. By Lemma 6.7, there exists a group subscheme $C$ of $H$, smooth over $S$, whose closed fiber is $H_s$ and such that $C_t=\Centr_{H_t}(T_t)$, where $T_t$ is a subtorus of $H_t$. Since $H$ is smooth over $S$, with connected fibers, one concludes, for reasons of dimension, that $C=H$. Keeping the notation of 6.7, one has $H=C=C_\ell$ for large $\ell$. Similarly consider $C'_\ell=\uCentr_G(M(\ell))$ (2.5), and let $C'$ be the stationary value of $C'_\ell$ for large $\ell$ (2.5 bis). The group scheme $C'$ contains $H$ and is such that $C'_t=\Centr_{G_t}(T_t)$ (6.8) and $C'_s=\Centr_{G_s}(T_s)$. The hypothesis on $H_s$ implies $\dim H_s=\dim C'_s$. Moreover, $\dim H_s=\dim H_t$ (since $H$ is smooth over $S$) and $\dim C'_t\leq\dim C'_s$ (Exp. VIB 4.1), so one has $\dim H_t=\dim C'_t$. But since $G_t$ is smooth and connected, $C'_t=\Centr_{G_t}(T_t)$ is smooth and connected, so finally one has $H_t=C'_t=\Centr_{G_t}(T_t)$.

% Source PDF p. 46.
\noindent b) Suppose that $H_s$ contains a Cartan subgroup $C_s$ of $G_s$, i.e. the centralizer in $(G_s)_{\mathrm{red}}$ of a maximal torus $T_s$ of $G_s$. By 6.7, there exists a group subscheme $C$ of $H$, smooth over $S$, with connected fibers, which lifts $C_s$. It plainly suffices to prove that $C_t$ contains a Cartan subgroup of $G_t$. Now by a), applied with $H=C$, $C_t$ is the centralizer in $G_t$ of a subtorus of $G_t$, hence contains a Cartan subgroup of $G_t$.

\noindent c) Suppose that $H_s$ is a parabolic subgroup of $G_s$. Let $N=\uNorm_G(H)$. By hypothesis, $H$ is equal to its connected normalizer, so $N_s$ is smooth, and consequently is equal to $\Norm_{(G_s)_{\mathrm{red}}}(H_s)=H_s$ (Exp. XII 8 bis). But then $N$ is flat over $S$. We shall see in Exp. XVI that under these conditions $G/N$ is representable. Since $H_s=N_s$ is a parabolic subgroup of $G$, $(G/N)_s$ is proper. Since $(G/N)$ has connected fibers and is flat over $S$, it follows from EGA III 5.5.1 that $(G/N)$ is proper over $S$. Thus $(G/N)_t=G_t/N_t$ is proper over $t$, and so is $G_t/H_t$, since $N_t/H_t$ is finite. It then suffices to apply the following lemma:
% typo?: the source says that H_s=N_s is a parabolic subgroup of G; kept as printed.
\begin{lemma}{6.9}\label{XV.6.9}
Let $k$ be a field, $G$ an algebraic $k$-group, $H$ a smooth connected algebraic subgroup of $G$, $N=\Norm_G H$; then if $\dim H=\dim N$ and $G/H$ is proper, $H$ is a parabolic subgroup of $G$.
\end{lemma}
Indeed, one may suppose $k$ algebraically closed and $G$ smooth and connected. The center $Z$ of $G$ is contained in $N$, and the hypothesis $\dim H=\dim N$ implies that $Z'=(Z)^0_{\mathrm{red}}$ is contained in $H$, so $G'=G/Z'$ is affine (Exp. XII 6.1). Replacing $G$ by $G'$ and $H$ by its image in $G'$, one reduces to the case where $G$ is affine (Exp. XIV 4.9), and Lemma 6.9 then follows from BIBLE 6 Th. 4. We have therefore proved Lemma 6.6.

To finish proving 6.4, we must prove that the $S$-preschemes representing $\mathscr{LC}$ (resp. $\mathscr{CT}$, resp. $\mathscr P$) are of finite presentation over $S$. This assertion is local on $S$, which allows one to suppose $S$ affine, then, since $G$ is of finite presentation, $S$ noetherian (EGA IV 8.9). We have just seen that the natural inclusions $\mathscr{CT}\to\mathscr L$ (resp. $\mathscr P\to\mathscr L$) are immersions, so the same holds for the inclusions $\mathscr{CT}\to\mathscr{LC}$ (resp. $\mathscr P\to\mathscr{LC}$), and consequently it suffices to prove that $\mathscr{LC}$ is representable by an $S$-prescheme of finite presentation.

Let us return to the notation introduced in 5.2. For every integer $n\geq0$, let $\mathscr L^n$ therefore be the subfunctor of $\mathscr L$ such that
\[
\mathscr L^n(S')=\{H\in\mathscr L(S')\text{ such that }\uNorm_{G_{S'}}(H)=\uNorm_{G_{S'}}H^{(n)}\}.
\]
The $S$-functor $\mathscr L^n$ is thus representable by an open subprescheme of $\mathscr L$, the sum of the $\mathscr L_r^n$. Each $\mathscr L_r^n$ is of finite presentation over $S$ (5.3) and is empty for $r>\sup_{s\in S}\dim G_s$ (which is a finite number, $S$ being quasi-compact), so $\mathscr L^n$ is of finite presentation over $S$. It suffices to prove that $\mathscr{LC}$ is contained in $\mathscr L^n$ for sufficiently large $n$.

For every point $s$ of $S$, let $d(s)$ be the smallest integer $n$ (finite or infinite) such that $\mathscr{LC}_{G_s}\subset\mathscr L^n_{G_s}$. It suffices to show that the function $d$ is bounded on $S$, for if $M$ is an upper bound, $\mathscr{LC}$ will be set-theoretically contained in $\mathscr L^M$, hence $\mathscr{LC}$ will be contained in $\mathscr L^M$, since the latter is an open subset of $\mathscr L$. An immediate constructibility argument reduces us to proving that if $S$ is noetherian and integral with generic point $t$, then the function $d$ is bounded on a neighborhood of $t$.

% Source PDF p. 47.
\par\noindent\emph{a) Reduction to the case where $G$ is smooth over $S$.}
Proceeding as in 6.2, one sees that, after changing $S$, one may suppose that $(G)_{\mathrm{red}}$ is a group prescheme smooth over $S$, which we shall denote by $G'$. Put $X=\mathscr{LC}_G$, $X'=\mathscr{LC}_{G'}$, and let $H$ (resp. $H'$) be the group subprescheme of $G_X$ (resp. $G'_{X'}$) universal for the functor $\mathscr{LC}_G$ (resp. $\mathscr{LC}_{G'}$). Since $H$ is smooth over $X$, $H\times_X(X_{\mathrm{red}})$ is reduced, hence contained in $G'_{X_{\mathrm{red}}}$, and it is an element of $\mathscr{LC}_{G'}(X_{\mathrm{red}})$, whence a canonical morphism $p:X_{\mathrm{red}}\to X'$. It is clear that $p$ is a monomorphism; let us show that $p$ is even an immersion. Let $N'$ be the normalizer of $H'$ in $G_{X'}$. The set of points $s$ of $X'$ such that the immersion $H'_s\to N'_s$ is an open immersion is an open subset $U$, and $H'|_U\to N|_U$ is an open immersion (Exp. VIB 2.5 and EGA IV 17.9.5). It follows from 5.1 iii) that $U$ is the largest open subset of $X'$ over which $H'$ is equal to its connected normalizer in $G_{X'}$, so $H'|_U\in\mathscr{LC}_G(U)$. One immediately deduces that $p$ is an isomorphism of $X_{\mathrm{red}}$ onto $U_{\mathrm{red}}$. If one can show that $\mathscr{LC}_G$ is of finite type when $G$ is smooth, $X'$ will be of finite type over $S$, hence $X_{\mathrm{red}}$ will be of finite type ($S$ is noetherian), and consequently will be contained in $\mathscr L_G^M$ for sufficiently large $M$; the same will hold for $X$.
% typo?: the source writes N|_U after defining N'; kept as printed.

\par\noindent\emph{b) Case where $S$ is the spectrum of an algebraically closed field $k$ of characteristic $p>0$ and $G$ is a smooth algebraic group over $k$.}
Instead of using the infinitesimal neighborhoods $H^{(n)}$ of a group subprescheme $H$ of $G$, we shall use the radicial subgroups ${}_{F^n}(H)$, kernels of the iterates of the Frobenius morphism in $H$ (Exp. VIIA 4), which is legitimate here, since ${}_{F^n}(H)$ is contained in the infinitesimal neighborhood of order $p^n$ of the unit section of $H$. If $T$ is a subtorus of $G$, ${}_{F^n}(T)={}_{p^n}(T)$, and it is immediate by duality that $\Centr_G({}_{p^n}(T))$ is equal to $\Centr_G(T)$ for sufficiently large $n$. One then has the following more precise proposition:
\begin{proposition}{6.10}\label{XV.6.10}
Let $k$ be a field of characteristic $p>0$, $G$ a smooth algebraic $k$-group, $T$ a maximal torus of $G_k$, $m$ the smallest integer such that
\[
\uCentr_G({}_{p^m}(T))^0=\uCentr_G(T)^0.
\]
Then, for every $k$-prescheme $S$ and every $H\in\mathscr{LC}(S)$, one has
\[
\uNorm_{G_S}(H)=\uNorm_{G_S}({}_{F^m}(H)).
\]
A fortiori, $\mathscr{LC}$ is contained in $\mathscr L^{p^m}$.
\end{proposition}
Since $H$ is smooth over $S$, of finite presentation over $S$, and has constant nilpotent rank (namely that of $G$), we shall see in the next paragraph (7.3) that the functor $\mathscr C_H$ of Cartan subgroups of $H$ is representable by an $S$-prescheme smooth over $S$ (the reader will check that the proof given of this property does not use the fact that $\mathscr{LC}_H$ is of finite type over $S$). It then follows from Exp. XIII 3.1 that one may consider the open subset $U$ of regular points of $H$.

Let $S'$ be an $S$-prescheme, $g$ an element of $G(S')$ normalizing ${}_{F^m}(H)_{S'}$. To prove that $g$ normalizes $H_{S'}$, it suffices to prove that $\operatorname{int}(g)U_{S'}$ is contained in $H_{S'}$; indeed, $H_{S'}\cap\operatorname{int}(g)H_{S'}$ will then contain an open subgroup of $\operatorname{int}(g)H$, hence will be equal to $\operatorname{int}(g)H$, since the latter has connected fibers. Replacing $S'$ by a suitable $S''$, then $S''$ by $S$, one reduces to proving that if $g\in G(S)$ normalizes ${}_{F^m}(H)$ and $u\in U(S)$, then $\operatorname{int}(g)u\in H(S)$.

% Source PDF p. 48.
Now let $C$ be the unique Cartan subgroup of $H_S$ which ``contains'' $u$ (Exp. XIII 3.2). It suffices for us to show that $C'=\operatorname{int}(g)C\subset H$. Since $H\in\mathscr{LC}_G(S)$, $C$ is also a Cartan subgroup of $G$, but the latter has maximal tori, so (Exp. XII 7.1 (a)) $C$ is the centralizer in $G_S^0$ of its unique maximal torus $T$. It follows from the definition of $m$ (and the fact that two maximal tori of $G$ are locally conjugate for fpqc (Exp. XII 7.1)) that one has
\[
C=(\Centr_G(T))^0=\uCentr_G({}_{p^m}(T))^0,
\]
whence, by conjugation by $g$,
\[
C'=\uCentr_G(\operatorname{int}(g)({}_{p^m}(T)))^0.
\]
But $\operatorname{int}(g)({}_{p^m}(T))$ is a subgroup of multiplicative type of $\operatorname{int}(g)({}_{F^m}(H))$, which is equal to ${}_{F^m}(H)$ ($g$ normalizes ${}_{F^m}(H)$), hence is contained in $H$. It then follows from Exp. XIII 2.1 (which is proved when the base is a field, but extends immediately to the case of an arbitrary base) that, for $C'$ to be contained in $H$, it suffices that one have
\[
\Lie C'\subset\Lie H.
\]
Now $\Lie C'=\operatorname{int}(g)(\Lie C)\subset\operatorname{int}(g)(\Lie H)$.

On the other hand, if $m\geq1$, which one is free to suppose, one has (using Exp. VIIA 4.1.2)
\[
\Lie H=\Lie({}_{F^m}(H))=\Lie(\operatorname{int}(g)({}_{F^m}(H)))=\operatorname{int}(g)(\Lie H).
\]
Thus $\Lie C'\subset\Lie H$, which completes the proof of 6.10.

We shall need another definition of the integer $m$ introduced in 6.10.
\begin{lemma}{6.11}\label{XV.6.11}
Let $G$ be a smooth algebraic group defined over an algebraically closed field $k$ of characteristic $p>0$, $T$ a maximal torus of $G$, $\mathfrak g=\Lie G$, and $R$ the family of nonzero characters of $T$ occurring in the representation of $T$ in $\mathfrak g$ induced by the adjoint representation of $G$. For every element $r\in R$, denote by $e_r$ the largest integer $n$ such that $p^n$ divides $r$ in the group of characters of $T$. Then $m=\sup_{r\in R}(e_r+1)$ if $R\ne\varnothing$, and $m=0$ otherwise.
\end{lemma}
Indeed, $\Centr_G(T)$ is smooth and contained in $\Centr_G({}_{p^m}(T))$, so
\[
\begin{aligned}
(\Centr_G T)^0=(\Centr_G({}_{p^m}(T)))^0
&\Longleftrightarrow \Lie(\Centr_G T)=\Lie(\Centr_G({}_{F^m}(T)))\\
&\Longleftrightarrow \mathfrak g^T=\mathfrak g^{{}_{p^m}(T)}\quad\text{(Exp. II 5.2.3)}.
\end{aligned}
\]
Now with the usual notation, one has
\[
\mathfrak g=\mathfrak g_0\oplus\coprod_{r\in R}\mathfrak g_r.
\]
Thus $\mathfrak g^T=\mathfrak g_0$ and $\mathfrak g^{({}_{p^m}(T))}=\mathfrak g_0+\coprod_{r\in R'}\mathfrak g_r$, where $R'$ is the part of $R$ formed by the characters of $T$ whose restriction to ${}_{p^m}(T)$ is zero. But a nonzero character $r$ of $T$ has zero restriction to ${}_{p^m}(T)$ if and only if $m\leq e_r$, whence the lemma.

% Source PDF p. 49.
\noindent c) Let us return to the proof of 6.4. We have reduced (according to point a) and the paragraph preceding it) to the case where $S$ is a noetherian integral scheme and $G$ is smooth over $S$. We must show that the function $d$ is bounded on a neighborhood of the generic point $t$ of $S$. After changing $S$, we may suppose that $G$ has a split maximal torus (6.2) $T=D_S(M)$. Let $\mathfrak g=\coprod_{\lambda\in M}\mathfrak g^\lambda$ be the decomposition of the Lie algebra of $G$ according to the characters of $T$, and let $R$ be the finite set of nonzero characters of $M$ such that $\mathfrak g^\lambda\ne0$. Let us then distinguish two cases:

\noindent\emph{1st case:} the point $t$, and consequently all points of $S$, have residue characteristic $p>0$. It is clear, in view of the foregoing, that the function $d$ is then bounded above by $p^m$, where $m$ is defined as in 6.11.

\noindent\emph{2nd case:} the point $t$ has residue characteristic zero. For every $\lambda\in R$, let $n_\lambda$ be the largest integer dividing $\lambda$ in the group $M$, and put $n=\prod n_\lambda$, $\lambda\in R$. For every prime number $q$ dividing $n$, denote by $S_q$ the closed subset of $S$ formed by the points of $S$ whose residue characteristic is equal to $q$, and let $U$ be the nonempty open subset (it contains $t$) complementary in $S$ to the union of the $S_q$. If now $s$ is a point of $U$, either $s$ has residue characteristic zero, and then $d(s)=1$ (5.6), or $s$ has residue characteristic $p>0$ which does not divide $n$, so the integer $m$ for the group $G_s$, defined in 6.11, is less than or equal to one. Moreover, it follows from Exp. VII\nde{Argument to be made explicit\ldots} that one has
\[
\uNorm_{G_{S'}}({}_{F^1}(H))=\uNorm_{G_{S'}}(\Lie H)=\uNorm_{G_{S'}}(H^{(1)}).
\]
Finally, it follows from 6.10 that if $H\in\mathscr{LC}_{G_s}(S')$, one has $\uNorm_{G_{S'}}(H)=\uNorm_{G_{S'}}(H^{(1)})$, and consequently $d(s)\leq1$, so $d$ is bounded by $1$ on $U$.

This completes the proof of 6.4.
\begin{corollary}{6.12}\label{XV.6.12}
Let $S$ be a prescheme and $G$ a group $S$-prescheme of finite presentation. Suppose that the nilpotent rank (resp. the dimension of the Borel subgroups) of the fibers of $G$ is a locally constant function on $S$; then the functor $\mathscr C$ of Cartan subgroups of $G$ (resp. the functor $\mathscr B$ of Borel subgroups of $G$) is representable by an $S$-prescheme of finite presentation over $S$.
\end{corollary}
Indeed, after restricting $S$, one may suppose that the nilpotent rank $\nu$ of the fibers is constant. But then it is clear that $\mathscr C$ is represented by the subprescheme, both open and closed, of the prescheme $X$ representing $\mathscr{LC}$ (6.4), over which the universal subgroup of $G_X$ for the functor $\mathscr{LC}$ has relative dimension $\nu$. The proof is analogous for the functor $\mathscr B$, taking account of the representability of the functor $\mathscr P$.

\sgasection{7}{Cartan subgroups of a smooth group}
\begin{proposition}{7.1}\label{XV.7.1}
Let $S$ be a prescheme, $G$ a group $S$-prescheme of finite presentation, smooth over $S$, with connected fibers.

% Source PDF p. 50.
\noindent i) Let $\mathscr{CT}$ be the $S$-functor $(\Sch/S)^\circ\to\Ens$ such that for every $S$-prescheme $S'$, one has

$\mathscr{CT}(S')=$ set of group subpreschemes $H$ of $G_{S'}$, smooth over $S'$, such that for every point $s'$ of $S'$, $H_{\bar s'}$ is the centralizer in $G_{\bar s'}$ of a subtorus of $G_{\bar s'}$.

Then $\mathscr{CT}$ is representable by an $S$-prescheme of finite presentation over $S$, smooth and quasi-projective over $S$.

\noindent ii) If $S$ is artinian and $H\in\mathscr{CT}(S)$, $H$ is the centralizer in $G$ of a subtorus of $G$. If $S$ is the spectrum of a henselian local ring and $H\in\mathscr{CT}(S)$, $H$ is the centralizer in $G$ of a subgroup of multiplicative type of $G$, étale over $S$.

\noindent iii) If $L$ is a group $S$-prescheme, smooth and of finite presentation over $S$, $i:L\to G$ a group $S$-monomorphism, and $H$ an element of $\mathscr{CT}(S)$, then $\uTransp_G(H,L)$ and $\uTransp_G^{\mathrm{str}}(H,L)$ (cf. Exp. VIII 6.5 e)) are representable by closed subpreschemes of $G$, smooth over $S$.

\noindent iv) If $H\in\mathscr{CT}(S)$, $H$ is closed in $G$, $N=\uNorm_G(H)$ is representable by a group subprescheme of $G$, closed and smooth over $S$; $N/H$ is representable by a group $S$-prescheme, separated over $S$, étale and of finite type over $S$; $G/N$ is representable by an $S$-prescheme smooth and quasi-projective over $S$.

\noindent v) Let $G'$ be a group $S$-prescheme of finite presentation over $S$, and $u:G\to G'$ a faithfully flat group $S$-morphism, so that $G'$ satisfies the same hypotheses as $G$ (Exp. VIB 9). Then if $H\in\mathscr{CT}_G(S)$, the image of $H$ under $u$ is representable by a group subprescheme $H'$ of $G'$ which is an element of $\mathscr{CT}_{G'}(S)$. Moreover, $H\to H'$ is faithfully flat, and if $H$ is the centralizer in $G$ of a torus $T$, $H'$ is the centralizer in $G'$ of the torus $T'=u(T)$.

\noindent vi) Under the conditions of v), consider the $S$-morphism
\[
\widetilde u:\mathscr{CT}_G\longrightarrow\mathscr{CT}_{G'},\qquad H\longmapsto H'=u(H).
\]
Then $\widetilde u$ is a faithfully flat quasi-compact morphism; if moreover $\Ker u$ is central, $\widetilde u$ is an isomorphism, the inverse isomorphism being $H'\mapsto u^{-1}(H')$.
\end{proposition}
\par\noindent\emph{Proof of ii).}
For the first assertion, one may suppose $S$ local artinian with closed point $s$. Let $H\in\mathscr{CT}(S)$ and $T_s$ the maximal central torus of $H_s$, which is already defined over $\kappa(s)$ (cf. 3.4). Since $H_s\in\mathscr{CT}(s)$, one has $H_s=\Centr_{G_s}T_s$. The group $H$ is smooth, so $T_s$ lifts (uniquely) to a subtorus $T$ of $H$, central in $H$ (Exp. IX 3.6 bis and Exp. IX 5.6). But then $H'=\uCentr_G T$ contains $H$ and has the same fiber as $H$; since $H$ is flat over $S$, one has $H=H'$ (Exp. VIB 2.5).

Suppose now that $S$ is the spectrum of a henselian local ring, which one may suppose noetherian by the usual reductions. Denote by $s$ the closed point of $S$, $T_s$ the maximal central torus of $H_s$, $q$ an integer invertible on $S$, $\ell$ a power of $q$, $T$ an $S$-torus whose closed fiber is isomorphic to $T_s$ (Exp. X 4.6). On the other hand, let ${}_\ell H$ be the ``kernel'' of raising to the $\ell$th power in $H$, and let $U_\ell$ be the largest open subset of ${}_\ell H$ which is étale over $S$. It then follows from 1.3 and the fact that $H$ is flat over $S$ that $U_\ell$ contains ${}_\ell T_s$. Since $S$ is henselian, there exists a unique $S$-morphism
\[
u_\ell:{}_\ell T\longrightarrow U_\ell
\]
which, on the closed fiber, induces the canonical immersion ${}_\ell T_s\to(U_\ell)_s$.
